\section{NN-S019: exponential positive arithmetic circuit cost for a local counting character}

\paragraph{Statement.}
Let $G=(V,E_G)$ be a simple graph on $n\ge2$ vertices of
maximum degree at most a fixed $d\ge1$. Suppose that, for every
$S\subseteq V$ with $|S|\le n/2$, its edge boundary satisfies
$|\partial S|\ge h|S|$, where $h>0$ is fixed.
Give each vertex two independent weights $y_{v,0},y_{v,1}$,
and each edge $e=uv$ four independent weights $z_{e,a,b}$.
Choose an orientation of each edge only to order these two bits.
Define the integral polynomial
\[
 P_G(y,z)=\sum_{x\in\{0,1\}^V}
      \prod_{v\in V}y_{v,x_v}
      \prod_{e=uv\in E_G}z_{e,x_u,x_v}.
\]
It has $2^n$ distinct monomials, all with coefficient one,
$2n+4|E_G|=O_d(n)$ variables and degree
$M=n+|E_G|=O_d(n)$. Every binary monotone arithmetic circuit
over $\mathbb R_{\ge0}$ computing it as a formal polynomial
has at least the following number of internal addition or multiplication gates:
\begin{equation}\label{eq:S019-lower}
 2^{\kappa n}\quad\hbox{gates},\qquad
            \kappa=\frac{h}{3(d+1)(d+h)}.
\end{equation}
The circuit model permits arbitrary nonnegative real constants,
addition and multiplication, and sharing of intermediate results.

In contrast, NN-C034 constructs a constant-free integral arithmetic
circuit for the very same polynomial of size
\[
 \operatorname{poly}_d(n)
 2^{O_d(n(\log\log n)^2/\sqrt{\log n})}=2^{o(n)},
\]
and polynomial formal degree. Thus, for these explicitly specified
local counting polynomials on any constant-expansion family,
the positive versus integral arithmetic size ratio is at least
$2^{\kappa n/2}$ for all sufficiently large $n$.
This is a consequence of E combined with the elementary
positive-rectangle argument proved below. It is not a new
general lower-bound technique or a historical-priority claim
for monotone graph polynomials.

\paragraph{The positive circuit gives a small balanced cover.}
There are $M$ weight blocks: one block for each vertex and one
for each edge. Every output monomial contains exactly one
variable from every block. Remove zero gates and gates which
do not contribute to the output. At each remaining gate $g$,
write $P_g$ for its nonzero polynomial. Fix a path from $g$
to the output and a nonzero monomial from every sibling at a
multiplication on this path. Their product is a monomial $c$.
Positivity implies that $pc$ is an output monomial for every
monomial $p$ of $P_g$. Consequently all monomials at $g$ have
the same block multidegree, with each entry zero or one.
Write $A_g$ for its set of occupied blocks and $r_g=|A_g|$.
This reasoning does not assume that the circuit is a formula.
In particular, a positive-degree gate cannot occur twice in
a contributing parse tree, since its block degrees would
then exceed the output degrees.

At an addition gate its two nonzero inputs have the same
block set. At a multiplication gate their block sets are
disjoint and their union is the output block set of that gate.
For a fixed output monomial, follow a contributing parse tree
from degree $M$ downwards. At sums choose a child contributing
the monomial. At products, while a child has degree greater
than $2M/3$, follow such a child. At the first product where
neither child has that degree, the parent still has degree
greater than $2M/3$. Choose the larger child: its degree is
strictly greater than $M/3$ and at most $2M/3$. Since
$h>0$ and $n\ge2$ imply $M\ge3$, this child has degree at
least two and is an internal gate. Thus every output monomial
passes through an internal gate
in the set
\[
             \mathcal B=\{g\text{ internal}:M/3\le r_g\le2M/3\}.
\]

For a fixed $g\in\mathcal B$, sum the products of sibling
polynomials over all paths from $g$ to the output, obtaining
a nonnegative context polynomial $C_g$. This polynomial is
used only for the proof; no efficient construction is needed.
Every term $P_gC_g$ has support contained in the output
support, and $C_g$ occupies exactly the complementary blocks.
The path reasoning just given proves the support cover
\begin{equation}\label{eq:S019-cover}
 \operatorname{supp}(P_G)\subseteq
       \bigcup_{g\in\mathcal B}\operatorname{supp}(P_gC_g).
\end{equation}
There are at most as many products as circuit gates. We need
neither coefficient equality in this cover nor disjointness
between its products.

\paragraph{A balanced product fixes many original bits.}
Consider any nonzero product $AB$ with nonnegative coefficients,
occupying complementary block sets $\mathcal A,\mathcal D$ of
sizes between $M/3$ and $2M/3$, whose support is contained in
$\operatorname{supp}(P_G)$. Let $B_0$ be the vertices whose
incident factor blocks meet both sides of this partition.
Here the unary vertex block is an incident factor, as is every
edge block containing the vertex. Put $b=|B_0|$.

For every $v\in B_0$, all monomials of $A$ and $B$ must
encode the same fixed bit at $v$. Indeed fix a monomial in
$B$. Each monomial in $A$ must be compatible with it, because
their product has positive coefficient and is a valid output
monomial. Since both sides contain a block involving $v$,
this fixes its bit in every monomial of $A$. Fixing one such
monomial of $A$ then fixes it in every monomial of $B$.
Internal consistency of each side also follows from extension
by any nonzero monomial on the other side. Thus the entire
product contains at most $2^{n-b}$ output monomials.

Let $V_A$ and $V_B$ be the vertices all of whose incident
factors belong respectively to $\mathcal A$ and $\mathcal D$.
They partition $V\setminus B_0$. Every factor in
$\mathcal A$ touches $V_A\cup B_0$. Since a vertex touches
at most $d+1$ factors,
\[
 |V_A|\ge\frac{M}{3(d+1)}-b,
 \qquad |V_B|\ge\frac{M}{3(d+1)}-b.
\]
There is no graph edge from $V_A$ to $V_B$: its factor
would have to belong to both sides. Take the smaller of
$V_A,V_B$, of size $a\le n/2$. All edges in its boundary
end in $B_0$, so expansion and the degree bound give
$ha\le db$. Combining with
$a\ge M/(3(d+1))-b$ yields
\[
 b\ge\frac{hM}{3(d+1)(d+h)}\ge\kappa n.
\]
This remains valid if the lower bound for $a$ before
rearrangement is negative. Therefore each product in
\eqref{eq:S019-cover} contains at most $2^{n-\kappa n}$
of the $2^n$ output monomials. If the circuit has $s$ gates,
then $2^n\le s2^{n-\kappa n}$, proving
\eqref{eq:S019-lower}.

\paragraph{The integral upper bound and nonempty family.}
The displayed definition of $P_G$ is exactly a universal
partition polynomial with $n$ Boolean variables, $n$ unary
factors and $|E_G|$ binary factors. Its density is at most
$1+d/2$. Thus every hypothesis of NN-C034 holds without
additional verifier wires, variable copying or a change
of the parameter. Its circuit output is the integral
polynomial above, so the comparison is in characteristic
zero before any reduction. Since the upper exponent is
$o(n)$, the claimed exponential ratio follows directly.

The expander hypothesis supplies an infinite explicit
family. One precise classical source is Reingold--Vadhan--
Wigderson, \emph{Entropy waves, the zig-zag graph product,
and new constant-degree expanders}, August~1, 2001 version,
Theorem~3.3, printed p.~11, and Section~5, pp.~16--18. In the published \emph{Annals of Mathematics}~155 (2002), 157--187 version, these are Theorem~3.3 and its proof on pp.~170--171, with the explicit starting graph in Lemma~5.1 and Proposition~5.2 on pp.~178--179. The readable published version is also arXiv:math/0406038v1. Its constant-degree
graphs have a normalized spectral bound $\lambda<1$
independent of their size. To pass from its allowed
multigraphs to simple graphs, remove loops and merge
parallel edges. For a set of at most half the vertices,
the spectral inequality applied to its centered indicator
gives at least $d(1-\lambda)|S|/2$ crossing edges with
multiplicity. Each unordered edge has multiplicity at
most $d$, so the resulting simple graph has at least
$(1-\lambda)|S|/2$ crossing edges and maximum degree
at most $d$. This gives the required fixed positive $h$.
The original primary version is
\url{https://www.math.ias.edu/~avi/PUBLICATIONS/MYPAPERS/SALIL/EXPANDERS/ANNALS/RVW01.pdf}. The published version is
\url{https://arxiv.org/pdf/math/0406038}.

\paragraph{Representation and Grothendieck translation.}
Let $T$ be the split torus over $\mathbb Z$ with independent
characters indexed by all the weights $y,z$. On the finite
etale Boolean scheme
$Y=(\operatorname{Spec}\mathbb Z\sqcup\operatorname{Spec}\mathbb Z)^V$,
give the free rank-one summand at section $x$ the character
equal to its displayed monomial. Its pushforward to
$\operatorname{Spec}\mathbb Z$ is an actual free $T$-module
$\mathcal V_G$ of rank $2^n$, with character $P_G$.
The dictionary between split-torus representations and
character gradings is SGA~3, Expos\'e~I,
4.7.2--4.7.3; direct sum and tensor product induce the
addition and multiplication of the representation
Grothendieck ring, as in SGA~6, Expos\'e~IV,
(1.5.2) and Definition~2.2.

A circuit building this class solely by direct sums,
tensor products, and nonnegative integer multiplicities
from the specified one-dimensional characters is a
monotone arithmetic circuit. It therefore needs
$2^{\kappa n}$ gates. E and NN-C034 instead give a
$2^{o(n)}$ expression in the representation ring using
virtual differences and shared subexpressions. The
resulting class is still the actual effective class
$[\mathcal V_G]$. This proves an exponential cost for
removing subtraction from these precise representation
expressions, despite their bounded-locality counting origin.

The lower bound concerns this arithmetic model. It does
not cover formulas with arbitrary Schur-functor primitives,
lattice-point-counting oracles, coefficient extraction,
hashing operations, arbitrary dual-character generators, or
the entire finite pushforward as a unit-cost primitive. It does not count intermediate
split-semiring values equal to the supported zero: by
NN-S009, a value alone does not retain cancellation
provenance. It gives no lower bound for unrestricted
Boolean circuits and no P versus NP or VP versus VNP
resolution. In particular, the upper bound remains
subexponential, not polynomial. Neither A nor M is used.


\paragraph{A polynomial monotone Boolean circuit for the same support.}
The formal-polynomial lower bound above persists alongside a much smaller
Boolean support circuit.  For each unary or binary factor tuple, let its
weight be an input bit recording whether that tuple is allowed.  Write
$a=(a_{v,b},a_{e,b,c})$ for all these bits.  Then
\[
 \operatorname{Supp}_G(a)=
 \bigvee_{x\in\{0,1\}^V}
       \left(\bigwedge_v a_{v,x_v}\ \wedge\!
                         \bigwedge_{e=uv}a_{e,x_u,x_v}\right)
\]
has a binary monotone Boolean circuit with $O(n^3+|E_G|)$ gates,
constructible in that time from $G$.  In particular, on the bounded-degree
expander family it has size $O(n^3)$, although every monotone arithmetic
circuit for the full formal counting polynomial needs $2^{\kappa n}$
gates.  No use of E is needed for this Boolean upper bound.

Here is the full construction and proof.  First use input bits
$f_{v,b},f_{e,b,c}$ recording that a tuple is forbidden.  A forbidden
unary tuple $(v,b)$ supplies the clause $x_v\ne b$.  A forbidden edge
tuple $(e=uv,b,c)$ supplies the clause
\[
                       (x_u\ne b)\vee(x_v\ne c).
\]
Thus consistency of the allowed tuples is exactly satisfiability of
this two-CNF.  It is useful that the set of possible clauses is fixed
by $G$: the input bits indicate which clauses are present.

Make an implication graph whose $2n$ vertices are literals
$L_{v,b}$ meaning $x_v=b$, with complement
$\overline{L_{v,b}}=L_{v,1-b}$.  A unary forbidden bit supplies
the edge $L_{v,b}\to L_{v,1-b}$.  A binary forbidden bit supplies
the two edges
\[
 L_{u,b}\longrightarrow L_{v,1-c},\qquad
 L_{v,c}\longrightarrow L_{u,1-b}.
\]
Each adjacency bit is the disjunction of the forbidden-tuple bits
that supply that edge.  The total number of such occurrences is
$2n+8|E_G|$, so all adjacency bits require $O(n+|E_G|)$ binary
OR gates, allowing constants for absent edges.  Every edge
$L\to M$ has its contrapositive edge
$\overline M\to\overline L$.

Enumerate the $2n$ literals as $1,\ldots,2n$.  Compute reachability
using the following Boolean circuits:
\[
 R^{(0)}_{ij}=A_{ij}\vee[i=j],\qquad
 R^{(k)}_{ij}=R^{(k-1)}_{ij}\vee
             \bigl(R^{(k-1)}_{ik}\wedge R^{(k-1)}_{kj}\bigr)
                     \quad(1\le k\le2n).
\]
Induction on $k$ shows that $R^{(k)}_{ij}$ records whether there is
a directed path from $i$ to $j$ whose internal vertices are among
$1,\ldots,k$.  The induction uses a simple path when there is a
path, so that a visit to $k$ splits it into the two required parts.
Consequently $R^{(2n)}$ is the reachability relation.  This circuit
uses $O(n^3)$ binary AND and OR gates.  Its inputs and operations
are all monotone in the forbidden-tuple bits.

The two-CNF is unsatisfiable precisely when
\[
 U_G(f)=\bigvee_{v\in V}
 \left(R^{(2n)}_{L_{v,0},L_{v,1}}\wedge
       R^{(2n)}_{L_{v,1},L_{v,0}}\right)=1.
\]
This is exactly the ordinary two-CNF strong-component criterion of
Aspvall--Plass--Tarjan, Theorem~1, printed p.~122, cited precisely
below.  We import that classical theorem rather than reprove it.

The entire circuit for $U_G$ uses only AND, OR, constants and the
forbidden inputs.  Form its Boolean dual by interchanging AND with
OR, interchanging constants zero with one, and replacing each
forbidden input label by the corresponding allowed input label.
Induction gate by gate, using De Morgan's laws, shows that the
dual circuit computes
\[
                            1-U_G(1-a)=\operatorname{Supp}_G(a).
\]
This is an actual monotone circuit on the allowed bits; no negation
gates are introduced at its inputs or internally.  Dualization
preserves the number of gates.  This proves the claimed bound.

The classical source for the implication graph and its exact
strong-component criterion is Aspvall--Plass--Tarjan,
\emph{A linear-time algorithm for testing the truth of certain
quantified Boolean formulas}, \emph{Information Processing Letters}
8 (1979), 121--123: the graph is defined on printed p.~121, and
Theorem~1 with its complete proof is on printed p.~122.
The primary article is available at
\url{https://mathweb.ucsd.edu/~sbuss/CourseWeb/Math268_2007WS/2SAT.pdf}.
Only the ordinary existential two-CNF theorem is used here; the
quantified extension in its Theorem~2 is not used.  The cubic circuit
construction and monotone dualization are explicit above, so no
circuit bound is being imported from the article's
linear-time random-access-machine algorithm.

\paragraph{Geometric form of the support contrast.}
For each fixed allowed-tuple bit string $a$, let
\[
 B_n=\mathbb Z[x_v:v\in V]/(x_v^2-x_v:v\in V)
                   \simeq\prod_{x\in\{0,1\}^V}\mathbb Z.
\]
For a vertex or edge factor $f$, let $p_{f,a}\in B_n$ be its
Boolean interpolation polynomial, taking value $a_{f,x|I_f}$
on section $x$.  Explicitly the unary interpolation is
$a_{v,0}(1-x_v)+a_{v,1}x_v$, and the binary interpolation is
the corresponding sum of four degree-two Boolean indicator
monomials multiplied by their allowed bits.  Each $p_{f,a}$
is an idempotent.  Put $e_a=\prod_f p_{f,a}$ and
\[
                 Y_a=\operatorname{Spec}(B_n/(1-e_a)).
\]
This is the finite etale scheme over $\mathbb Z$ consisting of
exactly the allowed assignment sections.  Equivalently one may
use all equations $p_{f,a}=1$ together with the Boolean equations:
for idempotents the ideals $(1-e_a)$ and $(1-p_{f,a}:f)$ agree.
One inclusion follows by expanding $1-\prod_f p_{f,a}$ as a
telescoping sum; the other follows from
$(1-p_{f,a})e_a=0$.

For every field $K$, the scheme $Y_a\times\operatorname{Spec}K$
is nonempty if and only if $\operatorname{Supp}_G(a)=1$.
Thus the nonemptiness of this explicitly given finite etale
subscheme is computable by the displayed $O(n^3)$ monotone
Boolean circuit in its allowed-tuple parameters.  Its rank is
$P_G(a)$ as an ordinary nonnegative integer, while the full
torus character before specialization is the exponentially
costly positive arithmetic object described above.  The map
from positive formal polynomials to Boolean support is a
semiring map with idempotent addition.  It is not an evaluation
of invertible torus characters at zero in the representation
ring, and it does not retain multiplicities.  This distinction
is exactly why the arithmetic lower bound and the Boolean
upper bound coexist.
